\documentclass[11pt]{article}
\usepackage{preamble}
\setlength{\parskip}{4pt}

\begin{document}

\daybanner{AP Statistics \textperiodcentered\ Unit 3 \textperiodcentered\ Topic 3.11 \textperiodcentered\ TEACHER KEY}{Reminders and Practice Completion}

\sectionbanner{CED TETHER}

{\small
\begin{itemize}[leftmargin=1.5em,itemsep=2pt,topsep=2pt]
  \item \textbf{ENDURING UNDERSTANDING: UNC-4} --- An interval of values should be used to estimate parameters, in order to account for uncertainty.
  \item \textbf{Skills: 4.B} --- Interpret statistical calculations and findings to assign meaning or assess a claim. | 4.D --- Justify a claim based on a confidence interval.
  \item UNC-4.M Interpret a confidence interval for a difference of proportions. [Skill 4.B]
  \item UNC-4.M.1 In repeated random sampling with the same sample size, approximately C\% of confidence intervals created will capture the difference in population proportions. UNC-4.M.2 Interpreting a confidence interval for difference between population proportions should include a reference to the sample taken and details about the population it represents.
  \item UNC-4.N Justify a claim based on a confidence interval for a difference of proportions. [Skill 4.D]
  \item UNC-4.N.1 A confidence interval for difference in population proportions provides an interval of values that may provide sufficient evidence to support a particular claim in context.
\end{itemize}
}
\textbf{Essential question:} Which claims about direction, effect size, causation, and population scope does a confidence interval actually support?\par
\textbf{DOK-3 skill:} 4.D \textperiodcentered\ \textbf{FRQ pattern:} {\small\texttt{confidence-\allowbreak{}interval-\allowbreak{}claim-\allowbreak{}effect-\allowbreak{}size-\allowbreak{}and-\allowbreak{}scope-\allowbreak{}adjudication}}\par
\textbf{DOK rationale:} Part (c) adjudicates competing claims by coordinating interval direction, a practical-effect threshold, long-run confidence meaning, random assignment, volunteer recruitment, and the resulting limits on causal and population conclusions.\par

\textbf{Self-paced bonus sheet.} Students take it when they choose and turn it in when they finish. Everything they need is printed on the sheet. Score part (c) only, E / P / I, by hand --- never AI-graded or auto-scored. (a) and (b) are the ladder up; use them to see where a thin (c) came from.\par

\textbf{Questions to ask}
\begin{itemize}[leftmargin=1.5em,itemsep=2pt,topsep=2pt]
  \item What does excluding zero establish, and what does it not establish?
  \item Which plausible interval values are smaller than a 10-point increase?
  \item What repeated process owns the 95 percent?
  \item Which design feature supports causation, and which missing feature blocks district-wide generalization?
\end{itemize}

\begin{callout}{\IconWarn\ LOOK FOR}{calloutred}
\begin{itemize}[leftmargin=1.5em,itemsep=2pt,topsep=2pt]
  \item Treating the point estimate as a guaranteed minimum effect.
  \item Assigning 95 percent probability to the fixed difference or to this one interval.
  \item Confusing random assignment with random sampling when stating population scope.
\end{itemize}
\end{callout}

\sectionbanner{SCORING PART (c) --- E / P / I}

\textbf{Expected elements}
\begin{itemize}[leftmargin=1.5em,itemsep=2pt,topsep=2pt]
  \item \textbf{judgment-1} (required) --- Uses the positive interval to support an increase and the lower endpoint 0.018 below 0.10 to reject a guaranteed 10-point minimum.
  \item \textbf{judgment-2} (required) --- Explains that 95 percent is long-run method performance, not a probability assigned to the fixed effect.
  \item \textbf{judgment-3} (required) --- Uses random assignment for the experimental causal comparison and volunteer recruitment to withhold district-wide generalization.
\end{itemize}

{\small\renewcommand{\arraystretch}{1.25}
\noindent\begin{tabular}{|p{0.5in}|p{5.9in}|}
\hline
\textbf{E} & Meets all three checks with correct contextual reasoning; equivalent wording is acceptable. \\ \hline
\textbf{P} & Shows at least one substantive correct check, but has an omission or error that prevents E. \\ \hline
\textbf{I} & Shows none of the three checks with substantive correct reasoning; a choice alone is insufficient. \\ \hline
\end{tabular}}

\clearpage
\focusbanner{THE PROBLEM --- annotated}

In a district experiment, 360 volunteers are randomly assigned: 180 get daily reminders and 180 use a standard platform. By Friday, 126 reminder and 105 standard students finish optional practice. For students like these volunteers, a 95\% interval for $p_R-p_S$ is $(0.018,0.215)$.\par\smallskip \textbf{Iris:} ``The positive interval and random assignment support a causal increase of 1.8 to 21.5 points for students like these volunteers.''\par \textbf{Cole:} ``There is a 95\% chance reminders raise completion for all district students by at least 10 points.''\par
\begin{center}
\textbf{Experiment counts}\par\smallskip
\begin{tabular}{|l|c|c|c|}
\hline
\textbf{Version} & \textbf{Complete} & \textbf{Other} & \textbf{Total} \\ \hline
\textbf{Reminder} & 126 & 54 & 180 \\ \hline
\textbf{Standard} & 105 & 75 & 180 \\ \hline
\end{tabular}
\end{center}
\begin{firsttakebox}{FIRST TAKE --- one sentence before you work the parts}
Which report would you share with classmates? Give one reason from the study.\par
\answer{Not graded. Cole joins familiar numbers into unsupported claims.}
\end{firsttakebox}

\begin{callout}{\IconBook\ MATCH EVERY CLAIM TO THE INTERVAL AND DESIGN (keep on the page)}{calloutgreen}
A confidence interval for $p_1-p_2$ gives plausible values for that population difference in the stated order.
If 0 is inside, no difference is plausible; if the whole interval is positive, the evidence supports $p_1>p_2$.
A claimed minimum effect $d$ is supported only when every interval value is at least $d$.
The confidence level is the long-run percent of intervals from repeated use of the method that capture the parameter,
not the probability that a fixed parameter or one computed interval is correct. Random assignment can support causation
for the experimental units; random sampling is what supports generalization to the sampled population.
\end{callout}

\dokbadge{1}~\textbf{(a)} State what $p_R-p_S$ represents, and identify whether $0$ and $0.10$ are contained in the reported interval.\par
\answer{Here $p_R-p_S$ is the reminder-minus-standard difference in the population proportions of students like these volunteers who would complete the optional practice under the two conditions. The value $0$ is not in $(0.018,0.215)$, while $0.10$ is in the interval.}

\dokbadge{2}~\textbf{(b)} Interpret the 95\% confidence interval in context. Then interpret the 95\% confidence level as the long-run performance of the interval procedure.\par
\answer{We are 95\% confident that, for students like these volunteers, the reminder completion proportion is 0.018 to 0.215 greater than the standard proportion. In repeated uses of this method with the same group sizes, about 95\% of constructed intervals would capture the represented population difference.}

\dokbadge{3}~\textbf{(c)} Choose the warranted conclusion. Use the positions of 0 and 0.10, distinguish random assignment from volunteer recruitment, and diagnose Cole's probability, minimum-effect, and population claims.\par
\answer{Iris is warranted. The entire interval is positive, so it supports a causal increase for these randomly assigned volunteers or similar students. But $0.10$ is inside and the lower endpoint is $0.018$, so an increase below 10 points remains plausible despite the point estimate $0.1167$. Also, 95\% is the method's long-run capture rate, not a probability about this fixed interval or effect. Volunteer recruitment is not random sampling, so the experiment does not represent every district student. Cole falsely assures a 10-point minimum, a 95\% probability, and district-wide scope.}

\end{document}
